% Daftar Pustaka (gaya APA edisi ke-7) -- contoh semua jenis sumber
\babtp{DAFTAR PUSTAKA}
Daftar pustaka berisi \textbf{semua} sumber yang dikutip di dalam teks, dan hanya sumber tersebut. Gunakan gaya APA edisi ke-7. Urutkan menurut abjad nama penulis pertama (atau judul bila tanpa penulis), tanpa nomor urut; baris kedua dan seterusnya menjorok ke dalam. Penulis dapat menggunakan bantuan Mendeley atau Zotero untuk menyusun sitasi.

Contoh di bawah mencakup jenis sumber yang umum dipakai: buku, buku terjemahan, bab dalam buku suntingan, artikel jurnal (dengan DOI), prosiding, skripsi atau tesis, laporan lembaga, lembar data (\textit{datasheet}), artikel daring, kode sumber di repositori, dan sumber tanpa penulis. Hapus contoh yang tidak dipakai dan ganti dengan sumber Anda sendiri.

\textbf{Sumber daring dicantumkan tanggal aksesnya.} Untuk laporan daring, artikel daring, dan repositori, akhiri entri dengan \textit{Diakses pada tanggal bulan tahun, dari alamat URL}, karena isinya dapat berubah. Tanggal akses adalah hari saat Anda membuka sumber tersebut. Artikel jurnal yang memiliki DOI cukup ditulis dengan DOI tanpa tanggal akses. Bila tahun terbit tidak diketahui, tulis (t.t.).

Bentuk sitasi di dalam teks untuk jenis sumber di atas:
\begin{itemize}
\item satu penulis: Andrews (2007) atau (Andrews, 2007);
\item dua penulis: (Xu \& Fischione, 2012);
\item tiga penulis atau lebih: (Capozzi et al., 2013);
\item lembaga: (Badan Pusat Statistik, 1997);
\item kutipan langsung dengan halaman: (Ghosh et al., 2010, hlm. 15);
\item tanpa penulis: (\textit{Panduan penulisan karya ilmiah}, 2020).
\end{itemize}

\begin{pustaka}
AlQahtani, S. A., \& Alhassany, M. (2013). Comparing different LTE scheduling schemes. \textit{2013 9th International Wireless Communications and Mobile Computing Conference (IWCMC)}, 264--269.

Andrews, M. (2007). A survey of scheduling theory in wireless data networks. \textit{Wireless Communications, 143}, 1--17.

Badan Pusat Statistik. (1997). \textit{Statistik Indonesia 1996}. Badan Pusat Statistik.

Capozzi, F., Piro, G., Grieco, L. A., Boggia, G., \& Camarda, P. (2013). Downlink packet scheduling in LTE cellular networks: Key design issues and a survey. \textit{IEEE Communications Surveys \& Tutorials, 15}(2), 678--700.

Espressif Systems. (2023). \textit{ESP32-WROOM-32 datasheet} (Versi 3.3) [Lembar data]. \url{https://www.espressif.com/en/support/documents/technical-documents}

Gaman, P. M., \& Sherrington, K. B. (1994). \textit{Ilmu pangan} (M. Gardjito, S. Naruki, Murdiati, \& Sardjono, Penerj.). Gadjah Mada University Press.

Ghosh, A., Zhang, J., Andrews, J. G., \& Muhamed, R. (2010). \textit{Fundamentals of LTE}. Prentice Hall.

Kementerian Komunikasi dan Digital. (2025). \textit{Judul laporan tahunan} [Laporan]. Diakses pada 21 Oktober 2026, dari \url{https://www.komdigi.go.id/contoh-laporan.pdf}

\textit{Panduan penulisan karya ilmiah}. (2020). Penerbit Universitas.

Penulis, A. (2020). Judul bab dalam buku suntingan. Dalam B. Editor \& C. Editor (Ed.), \textit{Judul buku suntingan} (hlm. 10--25). Penerbit.

Penulis, A. B. (2022). \textit{Judul skripsi atau tesis} [Skripsi, Universitas Muhammadiyah Buton]. Repositori Universitas Muhammadiyah Buton. \url{https://repository.example.ac.id/contoh}

Penulis, C. D. (2024). Judul artikel jurnal dengan DOI. \textit{Jurnal Teknologi dan Sistem Komputer, 12}(1), 45--58. \url{https://doi.org/10.xxxx/xxxxx}

Penulis, E. (2026, 5 Maret). Judul artikel berita atau blog. \textit{Nama Situs}. Diakses pada 21 Oktober 2026, dari \url{https://www.contoh.com/artikel}

Penulis, F. (2024). \textit{Nama repositori} [Kode sumber]. GitHub. Diakses pada 21 Oktober 2026, dari \url{https://github.com/pengguna/repositori}

Radiati, L. E. (2000). Penghambatan bakteri enteropatogen oleh ekstrak jahe. \textit{Jurnal Habitat, 12}(2), 81--91.

Xu, Y., \& Fischione, C. (2012). Real-time scheduling in LTE for smart grids. \textit{Proceedings of the 5th International Symposium on Communications, Control and Signal Processing}, 2--4.
\end{pustaka}
